\documentclass[11pt,a4paper]{article}
\usepackage[utf8]{inputenc}
\usepackage[italian]{babel}
\usepackage[T1]{fontenc}
\usepackage{cmap}
\usepackage{textcomp}
\usepackage{marvosym}
\usepackage{amssymb}
\usepackage[margin=1.5cm]{geometry}
\usepackage[table]{xcolor}
\usepackage{hyperref}
\usepackage{tcolorbox}
\usepackage{titlesec}
\usepackage{graphicx}
\usepackage{float}
\usepackage{booktabs}
\usepackage{enumitem}

% TPER palette - all defined with \definecolor (no inline mixing)
\definecolor{primaryColor}{HTML}{0054A6}
\definecolor{secondaryColor}{HTML}{4A5568}
\definecolor{backgroundColor}{HTML}{F7FAFC}
\definecolor{borderColor}{HTML}{E2E8F0}
\definecolor{successColor}{HTML}{38A169}
\definecolor{warningColor}{HTML}{D69E2E}
\definecolor{trainerColor}{HTML}{805AD5}
\definecolor{userColor}{HTML}{2B6CB0}
\definecolor{lightGray}{HTML}{EDF2F7}
\definecolor{midGray}{HTML}{A0AEC0}
\definecolor{darkGray}{HTML}{2D3748}
\definecolor{rogerColor}{HTML}{E53E3E}
\definecolor{channelBlue}{HTML}{3182CE}
\definecolor{channelGreen}{HTML}{38A169}
\definecolor{channelOrange}{HTML}{DD6B20}
\definecolor{channelTeal}{HTML}{319795}
\definecolor{cardBg}{HTML}{EBF4FF}

\colorlet{primaryLight}{primaryColor!10}
\colorlet{primaryMedium}{primaryColor!30}
\colorlet{trainerLight}{trainerColor!10}
\colorlet{successLight}{successColor!10}
\colorlet{warningLight}{warningColor!10}
\colorlet{cardBgLight}{cardBg!50}

\hypersetup{
    colorlinks=true,
    linkcolor=primaryColor,
    urlcolor=primaryColor,
}

\titleformat{\section}
  {\color{primaryColor}\normalfont\Large\bfseries}
  {}{0em}{}
  [\color{borderColor}\titlerule]

\titleformat{\subsection}
  {\color{primaryColor}\normalfont\large\bfseries}
  {}{0em}{}

% Screenshot style
\newcommand{\wfgraphic}[1]{%
  \includegraphics[width=\textwidth,keepaspectratio,frame]{#1}%
}

\begin{document}

\begin{tcolorbox}[colback=primaryColor,colframe=primaryColor,arc=3mm,halign=center,valign=center,height=2.5cm]
    \color{white}\LARGE\bfseries WIREFRAMES \\[0.2cm]
    \color{white!80}\large UX Design Project: Help People Help People --- TPER S.p.A. \\[0.1cm]
    \color{white!60}\normalsize Phase C --- Design Proposal --- Section 4.5
\end{tcolorbox}

\vspace{0.3cm}

\section{Introduzione}

I wireframe presentati in questa sezione illustrano le schermate fondamentali di \textbf{tper.it}, il portale informativo di TPER. A differenza di quanto ipotizzato in precedenza, \textbf{tper.it non consente acquisti diretti}: è una piattaforma informativa che supporta l'utente nella pianificazione del viaggio e nell'orientamento tra i canali di vendita fisici e digitali dell'ecosistema TPER.

\vspace{0.2cm}
\textbf{Canali dell'ecosistema TPER:}
\begin{itemize}
    \item \textbf{tper.it} --- Informativo: pianificazione percorsi, orari, tipi di biglietto, news, contatti
    \item \textbf{solweb.tper.it} --- Portale per l'acquisto di abbonamenti online
    \item \textbf{Roger App} --- App per biglietti digitali e pagamenti mobile
    \item \textbf{Punti Tper} --- Negozi fisici per tutti gli acquisti
    \item \textbf{Tabaccherie / PuntoLis} --- Punti di ricarica
    \item \textbf{Contactless a bordo} --- Pagamento per singola corsa con carta
\end{itemize}

\vspace{0.2cm}
\textbf{Convenzioni grafiche:}
\begin{itemize}
    \item Riquadri blu = elementi interattivi principali (CTA, navigazione)
    \item Riquadri grigi = contenuti informativi / aree secondarie
    \item Linea viola = elementi della Trainer-Mode
    \item Testo in \texttt{corsivo} = etichette descrittive
\end{itemize}

\newpage

% ================================================================
% WF1: HOMEPAGE tper.it (GUEST)
% ================================================================
\section{Wireframe 1: Homepage tper.it --- Utente Non Autenticato (Guest)}
\label{wf1}

\begin{figure}[H]
\centering
\wfgraphic{screenshot-wf1.png}
\caption{Homepage tper.it --- Vista utente non autenticato (Guest)}
\label{fig:wf1}
\end{figure}

\newpage

% ================================================================
% WF2: PIANIFICA (ROUTE PLANNER)
% ================================================================
\section{Wireframe 2: Pianifica --- Il Percorso (Core Function)}
\label{wf2}

\begin{figure}[H]
\centering
\wfgraphic{screenshot-wf2.png}
\caption{Pianifica --- Risultati di ricerca percorso con prezzi informativi}
\label{fig:wf2}
\end{figure}

\newpage

% ================================================================
% WF3: DOVE COMPRARE (CHANNEL GUIDE)
% ================================================================
\section{Wireframe 3: Dove Comprare --- Guida ai Canali}
\label{wf3}

\begin{figure}[H]
\centering
\wfgraphic{screenshot-wf3.png}
\caption{Dove Comprare --- Guida ai canali di vendita con tabella comparativa}
\label{fig:wf3}
\end{figure}

\newpage

% ================================================================
% WF4: BIGLIETTI (INFORMATIONAL)
% ================================================================
\section{Wireframe 4: Biglietti --- Tipologie e Prezzi}
\label{wf4}

\begin{figure}[H]
\centering
\wfgraphic{screenshot-wf4.png}
\caption{Biglietti --- Tipologie, prezzi e canali consigliati}
\label{fig:wf4}
\end{figure}

\newpage

% ================================================================
% WF5: AIUTO + TRAINER-MODE
% ================================================================
\section{Wireframe 5: Aiuto e Trainer-Mode --- Guida Passo-Passo}
\label{wf5}

\begin{figure}[H]
\centering
\wfgraphic{screenshot-wf5.png}
\caption{Aiuto e Trainer-Mode --- FAQ e guida passo-passo attivabile}
\label{fig:wf5}
\end{figure}

\vspace{0.5cm}

% ================================================================
% SUMMARY TABLE
% ================================================================
\section{Riepilogo dei Wireframe}

\begin{table}[H]
\centering
\small
\renewcommand{\arraystretch}{1.3}
\begin{tabular}{|p{2cm}|p{3cm}|p{4cm}|p{4.5cm}|}
\hline
\textbf{Wireframe} & \textbf{Schermata} & \textbf{Contenuto principale} & \textbf{Elementi chiave} \\\\[0.1cm]
\hline
WF1 & Homepage Guest & Route planner, 4 card informative (Pianifica, Biglietti, Abbonamenti, Aiuto), sidebar Dove Comprare & Nav tabs, hero search, canali in sidebar, news, contatti \\\\[0.1cm]
\hline
WF2 & Pianifica & Form di ricerca percorso, risultati con orari/durata/prezzo, toggle mappa/elenco & Core function, link Dove Comprare per ogni risultato, prezzi informativi \\\\[0.1cm]
\hline
WF3 & Dove Comprare & 5 schede canale (solweb, Roger, Punti Tper, Tabaccherie, Contactless), tabella comparativa & Card con CTA, confronto per esigenza, nessun acquisto diretto \\\\[0.1cm]
\hline
WF4 & Biglietti & Tipologie biglietto con prezzi e validita, badge consigliato, legenda canali & Informativo puro, icone canale per tipo, prezzi aggiornati \\\\[0.1cm]
\hline
WF5 & Aiuto + Trainer-Mode & FAQ per canale, sezione guide, toggle trainer-mode, guida 5 passi & Toggle ON/OFF, passo-passo visibile, canale indicato per ogni passo \\\\[0.1cm]
\hline
\end{tabular}
\caption{Riepilogo dei wireframe prodotti per tper.it (solo informativo)}
\label{tab:wfsum}
\end{table}

\vspace{0.3cm}

\begin{tcolorbox}[colback=primaryColor!5,colframe=primaryColor,arc=2mm]
\small
\textbf{Nota importante:} tper.it \`e un portale esclusivamente \textbf{informativo}. Non sono presenti flussi di acquisto, pagamento o transazione. Ogni schermata orienta l'utente verso il canale corretto (\texttt{solweb.tper.it}, Roger App, Punti Tper, Tabaccherie, Contactless a bordo) per completare l'operazione desiderata. Questa architettura rispetta la separazione reale dei canali nell'ecosistema TPER.
\end{tcolorbox}

\end{document}
